\chapter{Introduction}

\section{Objet du logiciel}

\logiciel{} est un logiciel desktop de modélisation et d'analyse de structures.
Il vise à offrir un environnement simple pour créer un modèle structurel,
définir les matériaux, sections, appuis, charges et combinaisons, puis lancer
une analyse et consulter les résultats.

Le logiciel est développé autour d'une interface graphique PySide6, d'une vue
3D interactive PyVista et d'un noyau multi-solveur. PyNite est le moteur de
calcul par défaut. OpenSeesPy reste disponible comme moteur avancé optionnel
lorsqu'il est installé sur la machine.

\section{Public visé}

Cette documentation s'adresse principalement :

\begin{itemize}
  \item aux ingénieurs structure qui souhaitent utiliser le logiciel ;
  \item aux développeurs qui veulent comprendre le fonctionnement général ;
  \item aux contributeurs qui enrichissent progressivement les fonctionnalités ;
  \item aux utilisateurs qui souhaitent produire une note de calcul ou un dossier
  de vérification à partir des résultats.
\end{itemize}

\section{Philosophie générale}

Le logiciel suit une logique métier classique :

\begin{enumerate}
  \item créer ou ouvrir un projet ;
  \item modéliser la géométrie : nœuds, barres, surfaces ;
  \item affecter matériaux, sections et conditions aux limites ;
  \item définir les cas de charge et les combinaisons ;
  \item lancer l'analyse ;
  \item lire les résultats sous forme de tableaux, diagrammes et vues graphiques.
\end{enumerate}

\begin{infobox}
Cette documentation est volontairement évolutive. Les sections marquées
\todoDoc{à compléter} servent de repères pour enrichir le manuel à mesure que le
logiciel avance.
\end{infobox}

\section{Unités internes}

Le système interne utilise principalement :

\begin{table}[H]
  \centering
  \begin{tabular}{lll}
    \toprule
    Grandeur & Unité interne & Exemple \\
    \midrule
    Longueur & \unite{m} & coordonnées des nœuds \\
    Force & \unite{kN} & charges ponctuelles \\
    Contrainte & \unite{kPa} & modules et résistances internes \\
    Moment & \unite{kN.m} & moments fléchissants \\
    \bottomrule
  \end{tabular}
  \caption{Unités internes principales}
\end{table}

\section{Statut du logiciel}

\logiciel{} est en développement actif. Les fonctionnalités de modélisation,
de calcul statique linéaire, de gestion des charges et de lecture des résultats
sont déjà présentes, mais le logiciel doit continuer à être validé sur des cas
de référence.

\begin{warningbox}
Les résultats produits par le logiciel ne constituent pas, à eux seuls, une note
de calcul certifiée. Toute utilisation professionnelle doit être vérifiée,
validée et signée par un ingénieur structure qualifié.
\end{warningbox}
